% assets/w06-rotation.svg — 회전행렬 R(psi) 가 무엇인가
%
%   bash _tools/tikz2svg.sh assets/src/w06-rotation.tex assets/w06-rotation.svg
%
% 이 그림이 말하는 것 넷
%   1) 선체 축은 NED 축을 psi 만큼 돌린 것이다. 그것이 전부다
%   2) 벡터는 하나인데 어느 축에서 읽느냐에 따라 숫자 두 쌍이 나온다
%   3) R 의 **열**은 선체 축의 단위벡터를 NED 에서 읽은 것이다 — 외울 것이 없다
%   4) 거꾸로 읽는 것은 전치다. R^-1 = R^T (직교행렬)
%
% 출처 — 제어공학특론 "1. 동역학 모델" (Fossen TTK4190 / QUT Robot Academy).
%        그 자료의 rot2 예제 수치(30도)를 그대로 쓴다
%
% 좌표 메모 (mm)
%   패널 원점  (1) 0   (2) 76   (3) 152
%   제목 y=60 · 설명 y=55 · 축 끝 y=36 · 원점 y=0 · 변수표 y=-20 · 가로줄 y=-36

\documentclass[border=6pt]{standalone}
% capstone-style.tex — 이 볼트의 TikZ 그림이 공유하는 스타일
%
% 쓰는 법:  % capstone-style.tex — 이 볼트의 TikZ 그림이 공유하는 스타일
%
% 쓰는 법:  % capstone-style.tex — 이 볼트의 TikZ 그림이 공유하는 스타일
%
% 쓰는 법:  \input{capstone-style}   (assets/src/ 안에서)
% 빌드:     bash _tools/tikz2svg.sh assets/src/w02-node-topic.tex assets/w02-node-topic.svg
%
% 왜 TikZ 인가
%   손으로 SVG 좌표를 쓰면 화살촉·선 굵기·글자 크기가 그림마다 미묘하게 달라진다.
%   그것이 "자료가 어설퍼 보인다"의 정체다. 스타일을 한 번만 정하고 상대좌표로
%   배치하면 좌표를 고쳐도 그림이 무너지지 않는다.
%   수식은 본문과 같은 TeX 글꼴로 나온다.
%
% 한글
%   ko.TeX(DVI 경로)를 쓴다. XeLaTeX 은 TikZ 그래픽을 PDF special 로 내보내는데
%   dvisvgm 이 그것을 받으려면 Ghostscript < 10.01 이 필요하다. 이 PC 에는
%   10.07 이 깔려 있어 그 경로가 막힌다. latex -> DVI -> dvisvgm 이 유일하게
%   동작하는 조합이다.

\usepackage[hangul]{kotex}
\usepackage{tikz}
\usepackage{amsmath}
\usepackage{xcolor}
\usetikzlibrary{arrows.meta, positioning, calc, fit, backgrounds, decorations.pathmorphing}

% ---- 색 --------------------------------------------------------------------
% 강조는 하나만. 나머지는 검정. 흑백 인쇄에서도 읽혀야 한다.
\definecolor{ink}{HTML}{1A1A1A}
\definecolor{soft}{HTML}{666666}
\definecolor{accent}{HTML}{7C3AED}
\definecolor{warn}{HTML}{B91C1C}

% 계층 색 — Simulink 모델의 블록 색과 같은 규칙을 쓴다
\definecolor{guid}{HTML}{1D4ED8}    % 유도
\definecolor{ctrl}{HTML}{EA580C}    % 제어
\definecolor{thr}{HTML}{D97706}     % 추진기
\definecolor{plant}{HTML}{16A34A}   % 운동모델
\definecolor{nav}{HTML}{7C3AED}     % 항법 · 신호처리
\definecolor{pathc}{HTML}{0D9488}   % 계획 경로 · 항적
\definecolor{errc}{HTML}{C2410C}    % 오차
\definecolor{flow}{HTML}{0369A1}    % 조류 · 바람

% ---- 선과 화살촉 -----------------------------------------------------------
\tikzset{
  sig/.style   = {ink, line width=0.5pt, -{Stealth[length=4pt, width=3pt]}},
  wire/.style  = {ink, line width=0.5pt},
  fb/.style    = {ink, line width=0.5pt, densely dashed,
                  -{Stealth[length=4pt, width=3pt]}},
  asig/.style  = {accent, line width=0.5pt, densely dashed,
                  -{Stealth[length=4pt, width=3pt]}},
  % 블록 — 흰 바탕, 직각, 검은 테두리. 색은 테두리에만 준다
  blk/.style   = {draw=ink, line width=0.55pt, fill=white,
                  minimum width=18mm, minimum height=8.5mm, inner sep=2pt, font=\small},
  gblk/.style  = {blk, draw=guid},
  cblk/.style  = {blk, draw=ctrl},
  tblk/.style  = {blk, draw=thr},
  pblk/.style  = {blk, draw=plant},
  nblk/.style  = {blk, draw=nav},
  % 합산점 · 분기점
  sum/.style   = {draw=ink, line width=0.55pt, fill=white, circle,
                  inner sep=0pt, minimum size=4.6mm, font=\scriptsize},
  dot/.style   = {ink, fill=ink, circle, inner sep=0pt, minimum size=1.5mm},
  % 글자
  sname/.style = {font=\small, inner sep=1.5pt},
  note/.style  = {font=\scriptsize, soft, inner sep=1.5pt},
  wnote/.style = {font=\scriptsize, warn, inner sep=1.5pt},
  anote/.style = {font=\scriptsize, accent, inner sep=1.5pt},
  axis/.style  = {ink, line width=0.5pt},
  tick/.style  = {font=\scriptsize, soft},
  % 묶음 상자 — 서브시스템 표시
  grp/.style   = {draw=soft, line width=0.4pt, densely dotted, rounded corners=1mm,
                  inner sep=3mm},
}
   (assets/src/ 안에서)
% 빌드:     bash _tools/tikz2svg.sh assets/src/w02-node-topic.tex assets/w02-node-topic.svg
%
% 왜 TikZ 인가
%   손으로 SVG 좌표를 쓰면 화살촉·선 굵기·글자 크기가 그림마다 미묘하게 달라진다.
%   그것이 "자료가 어설퍼 보인다"의 정체다. 스타일을 한 번만 정하고 상대좌표로
%   배치하면 좌표를 고쳐도 그림이 무너지지 않는다.
%   수식은 본문과 같은 TeX 글꼴로 나온다.
%
% 한글
%   ko.TeX(DVI 경로)를 쓴다. XeLaTeX 은 TikZ 그래픽을 PDF special 로 내보내는데
%   dvisvgm 이 그것을 받으려면 Ghostscript < 10.01 이 필요하다. 이 PC 에는
%   10.07 이 깔려 있어 그 경로가 막힌다. latex -> DVI -> dvisvgm 이 유일하게
%   동작하는 조합이다.

\usepackage[hangul]{kotex}
\usepackage{tikz}
\usepackage{amsmath}
\usepackage{xcolor}
\usetikzlibrary{arrows.meta, positioning, calc, fit, backgrounds, decorations.pathmorphing}

% ---- 색 --------------------------------------------------------------------
% 강조는 하나만. 나머지는 검정. 흑백 인쇄에서도 읽혀야 한다.
\definecolor{ink}{HTML}{1A1A1A}
\definecolor{soft}{HTML}{666666}
\definecolor{accent}{HTML}{7C3AED}
\definecolor{warn}{HTML}{B91C1C}

% 계층 색 — Simulink 모델의 블록 색과 같은 규칙을 쓴다
\definecolor{guid}{HTML}{1D4ED8}    % 유도
\definecolor{ctrl}{HTML}{EA580C}    % 제어
\definecolor{thr}{HTML}{D97706}     % 추진기
\definecolor{plant}{HTML}{16A34A}   % 운동모델
\definecolor{nav}{HTML}{7C3AED}     % 항법 · 신호처리
\definecolor{pathc}{HTML}{0D9488}   % 계획 경로 · 항적
\definecolor{errc}{HTML}{C2410C}    % 오차
\definecolor{flow}{HTML}{0369A1}    % 조류 · 바람

% ---- 선과 화살촉 -----------------------------------------------------------
\tikzset{
  sig/.style   = {ink, line width=0.5pt, -{Stealth[length=4pt, width=3pt]}},
  wire/.style  = {ink, line width=0.5pt},
  fb/.style    = {ink, line width=0.5pt, densely dashed,
                  -{Stealth[length=4pt, width=3pt]}},
  asig/.style  = {accent, line width=0.5pt, densely dashed,
                  -{Stealth[length=4pt, width=3pt]}},
  % 블록 — 흰 바탕, 직각, 검은 테두리. 색은 테두리에만 준다
  blk/.style   = {draw=ink, line width=0.55pt, fill=white,
                  minimum width=18mm, minimum height=8.5mm, inner sep=2pt, font=\small},
  gblk/.style  = {blk, draw=guid},
  cblk/.style  = {blk, draw=ctrl},
  tblk/.style  = {blk, draw=thr},
  pblk/.style  = {blk, draw=plant},
  nblk/.style  = {blk, draw=nav},
  % 합산점 · 분기점
  sum/.style   = {draw=ink, line width=0.55pt, fill=white, circle,
                  inner sep=0pt, minimum size=4.6mm, font=\scriptsize},
  dot/.style   = {ink, fill=ink, circle, inner sep=0pt, minimum size=1.5mm},
  % 글자
  sname/.style = {font=\small, inner sep=1.5pt},
  note/.style  = {font=\scriptsize, soft, inner sep=1.5pt},
  wnote/.style = {font=\scriptsize, warn, inner sep=1.5pt},
  anote/.style = {font=\scriptsize, accent, inner sep=1.5pt},
  axis/.style  = {ink, line width=0.5pt},
  tick/.style  = {font=\scriptsize, soft},
  % 묶음 상자 — 서브시스템 표시
  grp/.style   = {draw=soft, line width=0.4pt, densely dotted, rounded corners=1mm,
                  inner sep=3mm},
}
   (assets/src/ 안에서)
% 빌드:     bash _tools/tikz2svg.sh assets/src/w02-node-topic.tex assets/w02-node-topic.svg
%
% 왜 TikZ 인가
%   손으로 SVG 좌표를 쓰면 화살촉·선 굵기·글자 크기가 그림마다 미묘하게 달라진다.
%   그것이 "자료가 어설퍼 보인다"의 정체다. 스타일을 한 번만 정하고 상대좌표로
%   배치하면 좌표를 고쳐도 그림이 무너지지 않는다.
%   수식은 본문과 같은 TeX 글꼴로 나온다.
%
% 한글
%   ko.TeX(DVI 경로)를 쓴다. XeLaTeX 은 TikZ 그래픽을 PDF special 로 내보내는데
%   dvisvgm 이 그것을 받으려면 Ghostscript < 10.01 이 필요하다. 이 PC 에는
%   10.07 이 깔려 있어 그 경로가 막힌다. latex -> DVI -> dvisvgm 이 유일하게
%   동작하는 조합이다.

\usepackage[hangul]{kotex}
\usepackage{tikz}
\usepackage{amsmath}
\usepackage{xcolor}
\usetikzlibrary{arrows.meta, positioning, calc, fit, backgrounds, decorations.pathmorphing}

% ---- 색 --------------------------------------------------------------------
% 강조는 하나만. 나머지는 검정. 흑백 인쇄에서도 읽혀야 한다.
\definecolor{ink}{HTML}{1A1A1A}
\definecolor{soft}{HTML}{666666}
\definecolor{accent}{HTML}{7C3AED}
\definecolor{warn}{HTML}{B91C1C}

% 계층 색 — Simulink 모델의 블록 색과 같은 규칙을 쓴다
\definecolor{guid}{HTML}{1D4ED8}    % 유도
\definecolor{ctrl}{HTML}{EA580C}    % 제어
\definecolor{thr}{HTML}{D97706}     % 추진기
\definecolor{plant}{HTML}{16A34A}   % 운동모델
\definecolor{nav}{HTML}{7C3AED}     % 항법 · 신호처리
\definecolor{pathc}{HTML}{0D9488}   % 계획 경로 · 항적
\definecolor{errc}{HTML}{C2410C}    % 오차
\definecolor{flow}{HTML}{0369A1}    % 조류 · 바람

% ---- 선과 화살촉 -----------------------------------------------------------
\tikzset{
  sig/.style   = {ink, line width=0.5pt, -{Stealth[length=4pt, width=3pt]}},
  wire/.style  = {ink, line width=0.5pt},
  fb/.style    = {ink, line width=0.5pt, densely dashed,
                  -{Stealth[length=4pt, width=3pt]}},
  asig/.style  = {accent, line width=0.5pt, densely dashed,
                  -{Stealth[length=4pt, width=3pt]}},
  % 블록 — 흰 바탕, 직각, 검은 테두리. 색은 테두리에만 준다
  blk/.style   = {draw=ink, line width=0.55pt, fill=white,
                  minimum width=18mm, minimum height=8.5mm, inner sep=2pt, font=\small},
  gblk/.style  = {blk, draw=guid},
  cblk/.style  = {blk, draw=ctrl},
  tblk/.style  = {blk, draw=thr},
  pblk/.style  = {blk, draw=plant},
  nblk/.style  = {blk, draw=nav},
  % 합산점 · 분기점
  sum/.style   = {draw=ink, line width=0.55pt, fill=white, circle,
                  inner sep=0pt, minimum size=4.6mm, font=\scriptsize},
  dot/.style   = {ink, fill=ink, circle, inner sep=0pt, minimum size=1.5mm},
  % 글자
  sname/.style = {font=\small, inner sep=1.5pt},
  note/.style  = {font=\scriptsize, soft, inner sep=1.5pt},
  wnote/.style = {font=\scriptsize, warn, inner sep=1.5pt},
  anote/.style = {font=\scriptsize, accent, inner sep=1.5pt},
  axis/.style  = {ink, line width=0.5pt},
  tick/.style  = {font=\scriptsize, soft},
  % 묶음 상자 — 서브시스템 표시
  grp/.style   = {draw=soft, line width=0.4pt, densely dotted, rounded corners=1mm,
                  inner sep=3mm},
}


\begin{document}
\begin{tikzpicture}[x=1mm, y=1mm]

\pgfmathsetmacro{\psiD}{30}
\pgfmathsetmacro{\rot}{90-\psiD}
\pgfmathsetmacro{\SC}{12}                 % 1 m = 12 mm
\pgfmathsetmacro{\pbx}{2}                 % p^b = (2, 1) — 선체 축에서 읽은 값
\pgfmathsetmacro{\pby}{1}
\pgfmathsetmacro{\pnx}{round((\pbx*cos(\psiD) - \pby*sin(\psiD))*1000)/1000}
\pgfmathsetmacro{\pny}{round((\pbx*sin(\psiD) + \pby*cos(\psiD))*1000)/1000}
\pgfmathsetmacro{\SX}{\pny*\SC}           % 화면 x = 동 성분
\pgfmathsetmacro{\SY}{\pnx*\SC}           % 화면 y = 북 성분

% ===== 공통 — NED 축 =======================================================
\newcommand{\nedaxes}{%
  \draw[soft, line width=0.35pt, -{Stealth[length=3pt,width=2pt]}] (-3,0) -- (36,0);
  \draw[soft, line width=0.35pt, -{Stealth[length=3pt,width=2pt]}] (0,-3) -- (0,36);
  \node[font=\scriptsize, soft, anchor=north west] at (34,-0.6) {$y_n$ (동)};
  \node[font=\scriptsize, soft, anchor=south west] at (0.8,34)  {$x_n$ (북)};
}

% ===== 공통 — 선체와 선체 축 ===============================================
%   #1 선수축을 그릴 길이
\newcommand{\hull}[1]{%
  \begin{scope}[rotate=\rot]
    \draw[guid, line width=0.55pt, -{Stealth[length=3.4pt,width=2.4pt]}] (0,0) -- (#1,0);
    \node[font=\scriptsize, guid, anchor=west] at ({#1+1.2},0) {$x_b$ (선수)};
    \draw[guid, line width=0.55pt, -{Stealth[length=3.4pt,width=2.4pt]}] (0,0) -- (0,-15);
    \node[font=\scriptsize, guid, anchor=west] at (1.0,-15) {$y_b$ (우현)};
    \draw[ink, line width=0.5pt, fill=white]
         (-4.2,-2.7) -- (-4.2,2.7) -- (1.6,2.7) -- (5.2,0) -- (1.6,-2.7) -- cycle;
  \end{scope}}

% ===== 공통 — 선수각 호 =====================================================
\newcommand{\psiarc}[1]{%
  \draw[errc, line width=0.42pt, -{Stealth[length=3pt,width=2pt]}]
       (0,#1) arc (90:\rot:#1mm);
  \pgfmathsetmacro{\amid}{(90+\rot)/2}
  \pgfmathsetmacro{\arad}{#1+3.8}
  \node[font=\scriptsize, errc] at ({\arad*cos(\amid)},{\arad*sin(\amid)}) {$\psi$};}

% ===========================================================================
% (1) 선체 축은 NED 축을 psi 만큼 돌린 것이다
% ===========================================================================
\begin{scope}[shift={(0,0)}]
  \node[font=\small\bfseries, anchor=north west] at (-6,60) {(1) 축이 둘 — 하나는 땅, 하나는 배};
  \node[note, anchor=north west, text width=62mm] at (-6,55)
       {$x_n$ (북) $\cdot$ $y_n$ (동) 은 \textbf{땅에 붙어} 움직이지 않는다.
        $x_b$ (선수) $\cdot$ $y_b$ (우현) 는 \textbf{배에 붙어} 같이 돈다.
        둘 사이의 각이 선수각 $\psi$ 다};
  \nedaxes
  \hull{30}
  \psiarc{15}
  \node[anote, anchor=north west, text width=62mm] at (-6,-20)
       {선체 축은 NED 축을 $\psi$ 만큼 \textbf{돌린 것} — 회전행렬이 하는 일은 이것뿐이다};
\end{scope}

% ===========================================================================
% (2) 벡터 하나를 두 축에서 읽는다
% ===========================================================================
\begin{scope}[shift={(76,0)}]
  \node[font=\small\bfseries, anchor=north west] at (-6,60) {(2) 벡터는 하나, 숫자는 두 쌍};
  \node[note, anchor=north west, text width=62mm] at (-6,55)
       {보라색 화살표 $\mathbf{p}$ 는 \textbf{하나}다. 길이도 방향도 바뀌지 않는다.
        바뀌는 것은 \textbf{어느 자로 재느냐}뿐이다};
  \nedaxes
  \hull{30}

  % 벡터 p — 주인공
  \draw[accent, line width=1.8pt, -{Stealth[length=5pt,width=4pt]}] (0,0) -- (\SX,\SY);
  \node[font=\small, accent, anchor=south west] at ({\SX+0.6},{\SY+0.6}) {$\mathbf{p}$};

  % NED 성분 — 가로 · 세로 점선
  \draw[nav, line width=0.35pt, densely dashed] (\SX,\SY) -- (\SX,0);
  \draw[nav, line width=0.35pt, densely dashed] (\SX,\SY) -- (0,\SY);
  \draw[nav, line width=1.0pt, -{Stealth[length=3.6pt,width=2.6pt]}] (0,0) -- (\SX,0);
  \draw[nav, line width=1.0pt, -{Stealth[length=3.6pt,width=2.6pt]}] (0,0) -- (0,\SY);
  \node[font=\scriptsize, nav, anchor=north] at ({\SX*0.52},-1.2) {$y_n$};
  \node[font=\scriptsize, nav, anchor=east]  at (-1.2,{\SY*0.52}) {$x_n$};

  % 선체 성분 — 돌아간 축을 따라
  \begin{scope}[rotate=\rot]
    \draw[ctrl, line width=0.35pt, densely dashed]
         ({\pbx*\SC},{-\pby*\SC}) -- ({\pbx*\SC},0);
    \draw[ctrl, line width=0.35pt, densely dashed]
         ({\pbx*\SC},{-\pby*\SC}) -- (0,{-\pby*\SC});
    \draw[ctrl, line width=1.0pt, -{Stealth[length=3.6pt,width=2.6pt]}]
         (0,0) -- ({\pbx*\SC},0);
    \draw[ctrl, line width=1.0pt, -{Stealth[length=3.6pt,width=2.6pt]}]
         (0,0) -- (0,{-\pby*\SC});
    \node[font=\scriptsize, ctrl, anchor=south] at ({\pbx*\SC*0.5},0.8) {$x_b$};
    \node[font=\scriptsize, ctrl, anchor=north east]  at (-0.6,{-\pby*\SC*0.55}) {$y_b$};
  \end{scope}

  \node[font=\scriptsize, ctrl, anchor=north west] at (-6,-20)
       {선체 자로 재면\; $\mathbf{p}^{b} = (2.000,\; 1.000)$};
  \node[font=\scriptsize, nav, anchor=north west] at (-6,-26)
       {땅 자로 재면\;\; $\mathbf{p}^{n} = (\pgfmathprintnumber[fixed,precision=3]{\pnx},\;
        \pgfmathprintnumber[fixed,precision=3]{\pny})$};
\end{scope}

% ===========================================================================
% (3) R 의 열은 선체 축을 NED 에서 읽은 것이다
% ===========================================================================
\begin{scope}[shift={(152,0)}]
  \node[font=\small\bfseries, anchor=north west] at (-6,60) {(3) $R$ 의 열에는 외울 것이 없다};
  \node[note, anchor=north west, text width=62mm] at (-6,55)
       {선체 축의 \textbf{단위벡터}를 땅 자로 읽은 숫자 — 그것이 $R$ 의 열이다.
        첫 열이 $\mathbf{e}_{x_b}$, 둘째 열이 $\mathbf{e}_{y_b}$};
  \nedaxes
  \hull{30}

  % 선체 x 축 단위벡터  (NED 성분 = (cos psi, sin psi))
  \pgfmathsetmacro{\ux}{round(sin(\psiD)*1000)/1000}     % 동 성분
  \pgfmathsetmacro{\uy}{round(cos(\psiD)*1000)/1000}     % 북 성분
  \draw[thr, line width=1.6pt, -{Stealth[length=4.6pt,width=3.6pt]}]
       (0,0) -- ({\ux*22},{\uy*22});
  \draw[thr, line width=0.35pt, densely dashed] ({\ux*22},{\uy*22}) -- ({\ux*22},0);
  \draw[thr, line width=0.35pt, densely dashed] ({\ux*22},{\uy*22}) -- (0,{\uy*22});
  \node[font=\scriptsize, thr, anchor=east] at ({\ux*22-1.2},{\uy*22}) {$\mathbf{e}_{x_b}$};

  % 선체 y 축 단위벡터  (NED 성분 = (-sin psi, cos psi))
  \pgfmathsetmacro{\vx}{round(cos(\psiD)*1000)/1000}     % 동 성분
  \pgfmathsetmacro{\vy}{round(-sin(\psiD)*1000)/1000}    % 북 성분
  \draw[plant, line width=1.6pt, -{Stealth[length=4.6pt,width=3.6pt]}]
       (0,0) -- ({\vx*22},{\vy*22});
  \node[font=\scriptsize, plant, anchor=north] at ({\vx*22},{\vy*22-1.2}) {$\mathbf{e}_{y_b}$};

  \node[font=\scriptsize, thr, anchor=north west] at (-6,-20)
       {$\mathbf{e}_{x_b} = (\cos\psi,\;\sin\psi) = (0.866,\; 0.500)$};
  \node[font=\scriptsize, plant, anchor=north west] at (-6,-26)
       {$\mathbf{e}_{y_b} = (-\sin\psi,\;\cos\psi) = (-0.500,\; 0.866)$};
\end{scope}

% ===========================================================================
% 식 — 아래 띠
% ===========================================================================
\draw[wire, soft] (-6,-36) -- (216,-36);

\node[anchor=north west, font=\small] at (-6,-41)
     {$\underbrace{%
        \begin{bmatrix} x_n \\ y_n \end{bmatrix}}_{\text{땅 자}}
       =
       \underbrace{%
        \begin{bmatrix} \cos\psi & -\sin\psi \\ \sin\psi & \phantom{-}\cos\psi \end{bmatrix}}_{R(\psi)}
       \underbrace{%
        \begin{bmatrix} x_b \\ y_b \end{bmatrix}}_{\text{선체 자}}$};

\node[anchor=north west, font=\small] at (66,-41)
     {$\begin{bmatrix} x_b \\ y_b \end{bmatrix}
       = R(\psi)^{\mathsf T}\begin{bmatrix} x_n \\ y_n \end{bmatrix}
       = \begin{bmatrix} \phantom{-}\cos\psi & \sin\psi \\ -\sin\psi & \cos\psi \end{bmatrix}
         \begin{bmatrix} x_n \\ y_n \end{bmatrix}$};

\node[anchor=north west, font=\scriptsize, warn, text width=66mm] at (152,-40)
     {\textbf{거꾸로 읽기는 전치다.} $R^{-1} = R^{\mathsf T}$,\;
      $R R^{\mathsf T} = I$,\; $\det R = 1$.
      역행렬을 구할 일이 없다 — 행과 열을 바꾸기만 하면 된다};

\node[anchor=north west, font=\scriptsize, soft, text width=215mm] at (-6,-62)
     {수치 확인 ($\psi = 30^{\circ}$):\;
      $R = \begin{bmatrix} 0.866 & -0.500 \\ 0.500 & \phantom{-}0.866 \end{bmatrix}$,\quad
      두 열의 길이는 각각 1,\quad 두 열의 내적은 0 (서로 직각),\quad $\det R = 1$.\quad
      MATLAB 로는 \texttt{R = [cos(psi) -sin(psi); sin(psi) cos(psi)]} 한 줄이다};

\end{tikzpicture}
\end{document}
